% Pista ana-26s-q1 · Anàlisi
% Procedència: PAU Catalunya 2026 · Convocatòria de setembre · Sèrie 2 · Exercici 1
\documentclass[11pt,a4paper]{article}
\usepackage{pista}
\pistainfo{Catalunya 2026}{Convocatòria de setembre}{Sèrie 2}

\begin{document}

\section*{\textcolor{blauPista}{Exercici 1}}

\noindent Considereu la funció $f(x) = \dfrac{5x}{x^{2} + 1}$.

\subsection*{1.1. \normalfont Estudieu-ne el domini i els intervals de creixement i decreixement, així com els possibles màxims i mínims locals. \hfill\textnormal{[1~punt]}}

\pista{Per al domini, mira si el denominador $x^{2} + 1$ es pot anul·lar per algun valor real. Després deriva amb la regla del quocient i simplifica el numerador amb cura: és on és més fàcil equivocar-se de signe. Com que el denominador de $f'(x)$ és un quadrat (i no s'anul·la), el signe de $f'$ depèn només del numerador. Troba'n els zeros, fes una taula de signes i decideix on $f$ creix i on decreix; a cada canvi de monotonia hi haurà un màxim o un mínim local. Dona'n també l'ordenada, calculant-la a $f(x)$ (no a $f'(x)$).}

\subsection*{1.2. \normalfont Calculeu les equacions de totes les asímptotes d'aquesta funció (si en té) i dibuixeu un esbós de la seva gràfica. \hfill\textnormal{[0,75~punts]}}

\pista{Per a les asímptotes verticals, recorda què has vist del denominador en estudiar el domini. Per a les horitzontals, calcula $\lim\limits_{x\to+\infty} f(x)$ i $\lim\limits_{x\to-\infty} f(x)$ (compara els graus del numerador i del denominador, o divideix-los tots dos per $x^{2}$). Si n'hi ha d'horitzontals a totes dues bandes, ja no en pot haver d'obliqües. Per a l'esbós, combina-ho amb l'apartat anterior: el punt de tall amb els eixos, el signe de $f$ (pensa quin signe té el denominador) i els extrems locals. Recorda que una gràfica sí que pot tallar una asímptota horitzontal: només cal que s'hi acosti quan $x \to \pm\infty$.}

\subsection*{1.3. \normalfont Calculeu l'equació de la recta normal (és a dir, la recta perpendicular a la tangent) a la corba $y = f(x)$ en el punt d'abscissa $x = 2$. \hfill\textnormal{[0,75~punts]}}

\pista{Calcula primer el punt de la corba, $\bigl(2, f(2)\bigr)$, i el pendent de la tangent, $f'(2)$, aprofitant la derivada de l'apartat 1.1. La normal és perpendicular a la tangent, de manera que el seu pendent és $m = -\dfrac{1}{f'(2)}$: cal invertir el valor \emph{i} canviar-li el signe (l'error típic és fer només una de les dues coses, o fer servir directament $f'(2)$). Acaba amb l'equació punt-pendent $y - f(2) = m\,(x - 2)$.}

\end{document}
